\documentclass[10pt,a4paper]{amsart}
\usepackage[margin=19mm]{geometry}
\usepackage{amsmath,amssymb,mathtools}
\usepackage[scr=rsfs,cal=euler]{mathalfa}
\usepackage{xfrac}
\usepackage[unicode,hidelinks,bookmarks=false]{hyperref}
\newcommand{\ie}{i.e.\ }
\newcommand{\cf}{cf.\ }
\newcommand{\resp}{respectively\ }
\newcommand{\Subsection}{\subsection}
\providecommand{\nobreakdash}{\nobreak\hbox{-}}
\setlength{\emergencystretch}{2em}
\multlinegap0pt
\usepackage{amssymb}
\usepackage{tikz}
\newtheorem{theorem}{Theorem}
\newtheorem{lemma}{Lemma}
\newtheorem{claim}{Claim}
\title{$L_2$ Tur\'an Problems for Small Tournaments and Stability}
\author{Daniel I\v{l}kovi\v{c}}
\date{Source: arXiv:2609.05042v1 (CC BY 4.0)}
\begin{document}
\maketitle

\noindent\emph{Source and licence.} $L_2$ Tur\'an Problems for Small Tournaments and Stability. Version arXiv:2609.05042v1 (CC BY 4.0), distributed by arXiv under the Creative Commons Attribution 4.0 International licence, http://creativecommons.org/licenses/by/4.0/, as recorded in the arXiv licence metadata for this exact version and shown on the abstract page https://arxiv.org/abs/2609.05042v1. Full licence notice: \texttt{licenses/arxiv-2609.05042-CC-BY-4.0.txt}.\par\smallskip

\noindent\textit{Editorial scope.} Counted unit: the article's theorem thm:tt4 (source0.tex lines 182--185). The article's complete original proof of it is delivered with it (source0.tex lines 268--386, one \texttt{proof} environment of 119 lines, which the article places away from the statement). No mathematics is written, repaired or re-derived: the statement and the proof are the article's own lines, in the article's own notation, and no retained line is substituted or re-flowed.  Retained uncounted as the source-local context the proof needs: lines 234--243.  Nothing was omitted from any retained span, so the package carries no omission record.  No retained line contains a citation, so the wrapper carries no bibliography and no bibliography entry.  No retained line contains a figure environment, an image-inclusion command or drawing code, so no image is delivered and the package declares no asset dependency.  Editorial formatting only, in addition: an AMS \texttt{amsart} preamble that restates the article's own declarations and macros used by the retained lines, namely the environments \texttt{theorem}, \texttt{lemma}, \texttt{claim} on separate counters, together with the standard packages those lines need.  The retained environments print in this document as Theorem 1, Lemma 1, Claim 1 in order of appearance, whereas the article prints the same items with the numbering of its own full document.  The complete native source of the article as distributed in the arXiv e-print for this exact version is bundled unchanged as \texttt{sources/209409/source0.tex}.
\begin{lemma} \label{lem:tt3_free}
Let $S$ be a $TT_3$-free digraph on $n$ vertices.
Let $a$ be the maximum out-degree among all vertices in $S$.
Let $b = n - a$.
Then the sum of out-degrees $E(S)$ and the $L_2$ norm squared of the out-degree sequence $\|d^+_S\|_2^2$ in $S$ satisfy:
\begin{enumerate}
\item $E(S) \le 2ab$
\item $\|d^+_S\|_2^2 \le ab(a+b)$
\end{enumerate}
\end{lemma}

\begin{theorem}\label{thm:tt4}
For all $m \ge 1$, the maximum possible $L_2$ norm squared of the out-degree sequence $\text{ex}^+_{L_2}(m, TT_4)$ is exactly achieved by $T_3(m)$.
That is, $\text{ex}^+_{L_2}(m, TT_4) = \|d^+_{T_3(m)}\|_2^2$.
\end{theorem}

\begin{proof}[Proof of Theorem \ref{thm:tt4}]
To prove that the maximum possible $L_2$ norm squared of the out-degree sequence for any $TT_4$-free digraph on $m$ vertices is achieved by the complete directed 3-partite Tur\'an graph $T_3(m)$,
we will establish both a matching lower and upper bound.

\paragraph{Step 1: Lower Bound ($T_3(m)$ is $TT_4$-free)}
The complete directed 3-partite Tur\'an graph $T_3(m)$ partitions its $m$ vertices into three independent sets $V_1, V_2, V_3$ of sizes $n_1, n_2, n_3$ that are as equal as possible ($|n_i - n_j| \le 1$ and $n_1+n_2+n_3=m$).
Between any two vertices in different parts, edges exist in both directions (2-cycles).
No edges exist within any part.

A tournament on 4 vertices, $TT_4$, requires exactly one directed edge between every pair of its 4 vertices.
By the Pigeonhole Principle, any set of 4 vertices chosen from $T_3(m)$ must contain at least two vertices in the same part $V_i$.
Because $V_i$ is an independent set, there are zero edges between these two vertices.
Thus, they cannot form a tournament.
By consequence, $T_3(m)$ is $TT_4$-free, giving us a valid lower bound:
\[
\text{ex}^+_{L_2}(m, TT_4) \ge \|d^+_{T_3(m)}\|_2^2 = \sum_{i=1}^3 n_i(m-n_i)^2
\]

\paragraph{Step 2: Exact Upper Bound for $TT_4$-Free Digraphs}
Let $D$ be a $TT_4$-free digraph on $m$ vertices.
Let $\Delta$ be the maximum out-degree in $D$, achieved by some vertex v.
Define $S = N^+_D(v)$, meaning $|S| = \Delta$.

\begin{claim}
The induced subdigraph $D[S]$ is $TT_3$-free.
\end{claim}

\begin{proof}
If $D[S]$ contained a $TT_3$, the vertex $v$ (which points to all vertices in $S$) alongside this $TT_3$ would form a $TT_4$.
This contradicts $D$ being $TT_4$-free.
\end{proof}

We decompose the $L_2$ norm squared of the out-degree sequence of $D$, $\|d^+_D\|_2^2$, over $V(D) \setminus S$ and $S$:
\[
\|d^+_D\|_2^2 = \sum_{u \notin S} (d^+_D(u))^2 + \sum_{u \in S} (d^+_D(u))^2
\]

\begin{enumerate}
\item For $u \notin S$, the out-degree is bounded by the maximum $\Delta$.
With $m - \Delta$ such vertices:
\[
\sum_{u \notin S} (d^+_D(u))^2 \le (m - \Delta)\Delta^2
\]
\item For $u \in S$, its out-edges partition into $S$ and $V(D) \setminus S$.
The number of out-edges to $V(D) \setminus S$ is at most $|V(D) \setminus S| = m - \Delta$.
Thus, $d^+_D(u) \le d^+_S(u) + m - \Delta$.
Squaring this yields:
\[
(d^+_D(u))^2 \le (d^+_S(u))^2 + 2(m - \Delta)d^+_S(u) + (m - \Delta)^2
\]
Summing this inequality over all $\Delta$ vertices in $S$ gives:
\[
\sum_{u \in S} (d^+_D(u))^2 \le \|d^+_S\|_2^2 + 2(m - \Delta)E(S) + \Delta(m - \Delta)^2
\]
\end{enumerate}

Because $S$ is $TT_3$-free, we apply Lemma \ref{lem:tt3_free}.
Let $a$ be the maximum out-degree in $S$, and let $b = \Delta - a$.
Both $a$ and $b$ are non-negative integers.
We substitute $E(S) \le 2ab$ and $\|d^+_S\|_2^2 \le ab\Delta = ab(a+b)$:
\[
\sum_{u \in S} (d^+_D(u))^2 \le ab(a+b) + 4ab(m - \Delta) + \Delta(m - \Delta)^2
\]

Adding the bound for $V(D) \setminus S$, the total $L_2$ norm squared is:
\[
\|d^+_D\|_2^2 \le (m - \Delta)\Delta^2 + ab(a+b) + 4ab(m - \Delta) + \Delta(m - \Delta)^2
\]

Let $c = m - \Delta$.
Notice that $a, B, c$ are non-negative integers summing to $m$.
Substituting $\Delta = a+b$ and $c$:
\[
\|d^+_D\|_2^2 \le c(a+b)^2 + ab(a+b) + 4abc + (a+b)c^2
\]

Expanding this expression algebraically:
\[
c(a^2+2ab+b^2) + a^2b+ab^2 + 4abc + ac^2+bc^2 = a^2b + ab^2 + b^2c + bc^2 + c^2a + ca^2 + 6abc
\]
Notice that this uniquely factorizes into the exact $L_2$ norm squared for a 3-partite graph $F(a,b,c)$:
\[
F(a,b,c) = a(b+c)^2 + b(a+c)^2 + c(a+b)^2 = a^2b + ab^2 + b^2c + bc^2 + c^2a + ca^2 + 6abc
\]
Thus, we have shown that $\|d^+_D\|_2^2 \le F(a,b,c) $ for some partition $a+b+c=m$.

\paragraph{Step 3: Optimization Over Integer Partitions}
To maximize $F(a,b,c)$ over integers $a+b+c=m$, we rewrite it using elementary symmetric polynomials:
\[
F(a,b,c) = m(ab+bc+ca) + 3abc
\]

Suppose the partition is not as balanced as possible.
Then there exist two parts, say $a$ and $b$, such that $a \ge b + 2$.
Consider moving them closer by 1: $a' = a - 1$ and $b' = b + 1$ (leaving $c' = c$).
The net change is:
\[
\Delta F = F(a-1, b+1, c) - F(a,b,c)
\]
\[
= \big[m((a-1)(b+1) + (b+1)c + c(a-1)) + 3(a-1)(b+1)c\big] - \big[m(ab+bc+ca) + 3abc\big]
\]
\[
= m(a - b - 1) + 3c(a - b - 1) = (m + 3c)(a - b - 1)
\]

Since $a \ge b + 2$, we have $a - b - 1 \ge 1$.
Since $m \ge 1$ and $c \ge 0$, we have $m + 3c \ge 1$.
Therefore, $\Delta F \ge 1 > 0$ strictly.

Because balancing any two disparate parts strictly increases the $L_2$ norm squared, the unique global maximum is attained when no two parts differ by more than 1.
These balanced dimensions define the complete directed Tur\'an graph $T_3(m)$.
Thus:
\[
\|d^+_D\|_2^2 \le \max_{a+b+c=m} F(a,b,c) = \|d^+_{T_3(m)}\|_2^2
\]

Combined with our lower bound, the maximum is achieved by $T_3(m)$.
\end{proof}

\end{document}
